% Options for packages loaded elsewhere
\PassOptionsToPackage{unicode}{hyperref}
\PassOptionsToPackage{hyphens}{url}
\PassOptionsToPackage{dvipsnames,svgnames,x11names}{xcolor}
%
\documentclass[
  11pt,
]{article}
\usepackage{amsmath,amssymb}
\usepackage{setspace}
\usepackage{iftex}
\ifPDFTeX
  \usepackage[T1]{fontenc}
  \usepackage[utf8]{inputenc}
  \usepackage{textcomp} % provide euro and other symbols
\else % if luatex or xetex
  \usepackage{unicode-math} % this also loads fontspec
  \defaultfontfeatures{Scale=MatchLowercase}
  \defaultfontfeatures[\rmfamily]{Ligatures=TeX,Scale=1}
\fi
\usepackage{lmodern}
\ifPDFTeX\else
  % xetex/luatex font selection
  \setmainfont[]{Latin Modern Roman}
  \setmonofont[]{DejaVu Sans Mono}
\fi
% Use upquote if available, for straight quotes in verbatim environments
\IfFileExists{upquote.sty}{\usepackage{upquote}}{}
\IfFileExists{microtype.sty}{% use microtype if available
  \usepackage[]{microtype}
  \UseMicrotypeSet[protrusion]{basicmath} % disable protrusion for tt fonts
}{}
\makeatletter
\@ifundefined{KOMAClassName}{% if non-KOMA class
  \IfFileExists{parskip.sty}{%
    \usepackage{parskip}
  }{% else
    \setlength{\parindent}{0pt}
    \setlength{\parskip}{6pt plus 2pt minus 1pt}}
}{% if KOMA class
  \KOMAoptions{parskip=half}}
\makeatother
\usepackage{xcolor}
\usepackage[margin=2.2cm]{geometry}
\usepackage{longtable,booktabs,array}
\usepackage{calc} % for calculating minipage widths
% Correct order of tables after \paragraph or \subparagraph
\usepackage{etoolbox}
\makeatletter
\patchcmd\longtable{\par}{\if@noskipsec\mbox{}\fi\par}{}{}
\makeatother
% Allow footnotes in longtable head/foot
\IfFileExists{footnotehyper.sty}{\usepackage{footnotehyper}}{\usepackage{footnote}}
\makesavenoteenv{longtable}
\setlength{\emergencystretch}{3em} % prevent overfull lines
\providecommand{\tightlist}{%
  \setlength{\itemsep}{0pt}\setlength{\parskip}{0pt}}
\setcounter{secnumdepth}{5}
\ifLuaTeX
\usepackage[bidi=basic]{babel}
\else
\usepackage[bidi=default]{babel}
\fi
\babelprovide[main,import]{french}
\ifPDFTeX
\else
\babelfont{rm}[]{Latin Modern Roman}
\fi
% get rid of language-specific shorthands (see #6817):
\let\LanguageShortHands\languageshorthands
\def\languageshorthands#1{}
\usepackage{float}
\usepackage{fvextra}
\usepackage{microtype}
\usepackage{etoolbox}
\usepackage{fancyhdr}
\AtBeginEnvironment{CSLReferences}{\small\setstretch{1.0}\interlinepenalty=10000}
\pagestyle{fancy}
\fancyhf{}
\fancyfoot[C]{\small Article 1 / V11 / Author-review copy \quad \thepage}
\renewcommand{\headrulewidth}{0pt}
\setlength{\emergencystretch}{3em}
\setlength{\tabcolsep}{4pt}
\floatplacement{figure}{H}
\ifLuaTeX
  \usepackage{selnolig}  % disable illegal ligatures
\fi
\IfFileExists{bookmark.sty}{\usepackage{bookmark}}{\usepackage{hyperref}}
\IfFileExists{xurl.sty}{\usepackage{xurl}}{} % add URL line breaks if available
\urlstyle{same}
\hypersetup{
  pdftitle={Article 1 - dossier de validation et de soumission},
  pdfauthor={Youssef EL MOUTEE - document de travail},
  pdflang={fr-FR},
  colorlinks=true,
  linkcolor={Maroon},
  filecolor={Maroon},
  citecolor={Blue},
  urlcolor={teal},
  pdfcreator={LaTeX via pandoc}}

\title{Article 1 - dossier de validation et de soumission}
\author{Youssef EL MOUTEE - document de travail}
\date{6 septembre 2026 - V11}

\begin{document}
\maketitle

\setstretch{1.12}
\hypertarget{sujet-et-plan-conservuxe9s}{%
\section{Sujet et plan conservés}\label{sujet-et-plan-conservuxe9s}}

Le sujet de thèse et le plan des articles restent ceux déjà engagés.
L'Article 1 conserve son titre original : \textbf{Reference-Set
Decoupled Co-Adaptive Training for Physics-Informed Neural Networks}.
Les corrections portent sur la précision des affirmations, la
bibliographie et la présentation du dossier.

\hypertarget{de-quelle-soumission-parle-t-on}{%
\section{De quelle soumission parle-t-on
?}\label{de-quelle-soumission-parle-t-on}}

Il s'agit de proposer \textbf{un article scientifique} à une revue. Cela
se distingue de l'inscription au doctorat, du dépôt de la thèse et de la
soutenance. Le dépôt sur GitHub conserve le travail ; il ne soumet pas
l'article à une revue.

L'auteur correspondant, désigné par l'équipe d'auteurs, ouvre le dossier
sur la plateforme de la revue, renseigne les auteurs et joint manuscrit,
figures, annexes et déclarations. Youssef peut tenir ce rôle s'il est
désigné ; l'encadrant peut aussi le tenir. La qualité d'encadrant ne
suffit pas, à elle seule, à fixer la liste ou l'ordre des auteurs.

Aucune date n'a été fixée. Le rappel prévu le 7 septembre est un point
d'avancement, pas une échéance de dépôt. La soumission pourra avoir lieu
après validation scientifique, informations d'auteurs complètes et
approbation finale. Aucune revue n'a reçu ce manuscrit dans ce travail.

Après le dépôt, la rédaction peut rejeter l'article d'emblée ou le
transmettre à des évaluateurs. Ceux-ci peuvent demander des révisions.
L'auteur correspondant transmet ensuite une version révisée et une
réponse point par point. L'acceptation éventuelle précède les épreuves
et la publication. Les délais dépendent de la revue et des révisions ;
aucune durée n'est promise.

\hypertarget{revue-retenue-pour-pruxe9parer-le-dossier}{%
\section{Revue retenue pour préparer le
dossier}\label{revue-retenue-pour-pruxe9parer-le-dossier}}

\textbf{Cible de travail : Journal of Computational Science, Elsevier,
ISSN 1877-7503.} Mon appréciation est que l'étude comparative, ses
contrôles numériques et sa reproductibilité correspondent au périmètre
de recherche par simulation affiché par la revue. Cette adéquation ne
préjuge pas de l'acceptation.

\begin{longtable}[]{@{}
  >{\raggedright\arraybackslash}p{(\columnwidth - 2\tabcolsep) * \real{0.5000}}
  >{\raggedright\arraybackslash}p{(\columnwidth - 2\tabcolsep) * \real{0.5000}}@{}}
\toprule\noalign{}
\begin{minipage}[b]{\linewidth}\raggedright
Critère
\end{minipage} & \begin{minipage}[b]{\linewidth}\raggedright
Décision de préparation
\end{minipage} \\
\midrule\noalign{}
\endhead
\bottomrule\noalign{}
\endlastfoot
Type de travail & Article de recherche numérique ; évaluation de
ReCoA-PINN et de mécanismes concurrents \\
Format & Source LaTeX modifiable et PDF de lecture ; tableaux et figures
numérotés \\
Évaluation & Le guide officiel indexé indique une évaluation en simple
anonymat \\
Highlights & Cinq points, chacun inférieur ou égal à 85 caractères \\
Accès & Voie par abonnement envisagée : le site indique l'absence de
frais de publication pour les auteurs sur cette voie \\
Accès ouvert & Option distincte ; aucun coût ni engagement n'est accepté
dans cette préparation \\
Décision finale & À valider avec l'encadrant et les auteurs avant
dépôt \\
\end{longtable}

Alternative : \textbf{Journal of Computational Physics} si les auteurs
jugent la contribution méthodologique suffisamment forte. La revue cible
reste unique lors d'une soumission ; aucun dépôt simultané n'est
préparé.

Sources consultées le 6 septembre 2026 :
\href{https://www.sciencedirect.com/journal/journal-of-computational-science}{périmètre
de Journal of Computational Science},
\href{https://www.sciencedirect.com/journal/journal-of-computational-science/publish/guide-for-authors}{guide
des auteurs},
\href{https://www.sciencedirect.com/journal/journal-of-computational-science/about/insights}{options
de publication}. Le guide complet a refusé l'ouverture directe ; les
exigences ci-dessus proviennent des extraits officiels indexés.
L'ensemble des champs de la plateforme devra être revérifié au moment du
dépôt. Le résumé est volontairement limité à moins de 250 mots comme
choix de préparation, sans présenter ce seuil comme vérifié pour cette
revue.

\hypertarget{auteurs-et-duxe9clarations-uxe0-confirmer}{%
\section{Auteurs et déclarations à
confirmer}\label{auteurs-et-duxe9clarations-uxe0-confirmer}}

\begin{longtable}[]{@{}
  >{\raggedright\arraybackslash}p{(\columnwidth - 2\tabcolsep) * \real{0.5000}}
  >{\raggedright\arraybackslash}p{(\columnwidth - 2\tabcolsep) * \real{0.5000}}@{}}
\toprule\noalign{}
\begin{minipage}[b]{\linewidth}\raggedright
Élément
\end{minipage} & \begin{minipage}[b]{\linewidth}\raggedright
Information disponible / validation requise
\end{minipage} \\
\midrule\noalign{}
\endhead
\bottomrule\noalign{}
\endlastfoot
Auteur identifié dans le projet & Youssef EL MOUTEE \\
Encadrant identifié par Youssef & Professeur Adil Echchelh \\
Coauteurs et ordre & Non confirmés ; ne pas ajouter automatiquement
l'encadrant \\
Affiliation de Youssef & À confirmer ; aucune affiliation universitaire
déduite \\
Auteur correspondant et adresse & À désigner et à confirmer \\
ORCID & À renseigner si disponible ; aucun identifiant inventé \\
Contributions CRediT & À attribuer selon les contributions effectivement
réalisées \\
Financement & À confirmer ; « aucun financement » n'est pas présumé \\
Conflits d'intérêts & Déclaration de chaque auteur à recueillir \\
Validation du contenu et du dépôt & Accord explicite de tous les auteurs
à obtenir \\
\end{longtable}

Le manuscrit de lecture identifie Youssef comme auteur de travail et
signale les métadonnées non finalisées. Il ne contient ni signature, ni
accord attribué à un coauteur.

\hypertarget{duxe9claration-dassistance-ia-proposuxe9e}{%
\subsection{Déclaration d'assistance IA
proposée}\label{duxe9claration-dassistance-ia-proposuxe9e}}

Le projet a utilisé ChatGPT/Codex d'OpenAI pour assister la préparation
du code, les scripts d'expérimentation et d'analyse, la recherche
bibliographique et la rédaction. Les modèles PINNs entraînés dans
l'étude sont distincts de cet assistant. Les résultats numériques
proviennent des exécutions et fichiers conservés. Les auteurs doivent
vérifier les méthodes, références, résultats et formulations et assumer
la responsabilité de la version finale avant soumission. La version
exacte de l'assistant doit être complétée à partir des informations
disponibles ; elle n'est pas inventée ici.

La
\href{https://www.elsevier.com/about/policies-and-standards/generative-ai-policies-for-journals}{politique
Elsevier sur l'IA générative} demande une déclaration de l'assistance
utilisée pour préparer le manuscrit et une description dans les méthodes
lorsqu'elle intervient dans la recherche. Le texte final de
responsabilité ne sera validé qu'après la relecture effective des
auteurs.

\hypertarget{ruxe9sultat-de-la-vuxe9rification-sur-une-autre-machine}{%
\section{Résultat de la vérification sur une autre
machine}\label{ruxe9sultat-de-la-vuxe9rification-sur-une-autre-machine}}

L'exécution du 6 septembre 2026 sur un serveur Ubuntu GitHub indépendant
a restauré les trois campagnes et passé \textbf{55 tests automatisés}.
Les \textbf{560 scores} et les \textbf{96 contrastes statistiques} sont
reproduits aux tolérances fixées avant l'exécution. Les \textbf{12
réentraînements} sont terminés ; leurs erreurs L2 et maximales
respectent toutes les tolérances numériques prévues. L'écart absolu le
plus élevé vaut \(9{,}53\times10^{-9}\) en L2 et \(1{,}40\times10^{-8}\)
en erreur maximale.

Le contrôle strict comprenant l'identité binaire des données initiales
reste à \textbf{8 réussites sur 12}. Les modèles initiaux sont
identiques dans les douze cas ; les empreintes des jeux d'entraînement
et de référence diffèrent pour les quatre essais bruités.
L'environnement original utilisait PyTorch 2.14.0+cpu ; l'environnement
public utilise 2.8.0+cpu. Les tableaux de données initiaux originaux
n'ayant pas été archivés, les octets différents ne peuvent pas être
localisés rétrospectivement. Un effet numérique de version ou
d'environnement est plausible, sans preuve permettant de l'attribuer à
une opération précise.

La conclusion retenue est donc une \textbf{reproduction numérique
réussie, avec une limite documentée sur l'identité exacte des entrées}.
Le fichier d'échec initial et les seuils restent inchangés. On ne
remplace pas le résultat par un succès artificiel. Si une identité bit à
bit est exigée, il faudra une campagne prospective archivant directement
les tenseurs initiaux dans un environnement figé ; ce résultat n'est pas
revendiqué dans le manuscrit.

Preuves :
\href{https://github.com/ucef90/Doctorat/actions/runs/34021530276}{exécution
indépendante}, résumés et CSV conservés dans
\texttt{submission/v11/reproduction/}.

\hypertarget{contribution-proposuxe9e-uxe0-la-validation-scientifique}{%
\section{Contribution proposée à la validation
scientifique}\label{contribution-proposuxe9e-uxe0-la-validation-scientifique}}

La question de recherche reste celle du projet : l'utilisation d'un
ensemble de référence fixe pour contrôler l'apprentissage coadaptatif
des PINNs. Les expériences permettent une étude comparative contrôlée,
mais ne démontrent ni une supériorité générale de M6, ni un gain de
temps. Le titre original est conservé, avec un résumé qui annonce ces
limites.

La relecture interne a précisé l'identité d'importance M7, la dépendance
des trajectoires adaptatives, le rôle du jeu de référence, l'écrêtage
suivi de renormalisation, les équations manufacturées et les familles
statistiques. La bibliographie a été actualisée à partir des
publications ou prépublications primaires, avec douze références dans le
manuscrit. La nouveauté suffisante pour publication et l'interprétation
finale doivent encore être examinées par l'encadrant et un lecteur
scientifique humain.

\hypertarget{avancement-des-cinq-demandes}{%
\section{Avancement des cinq
demandes}\label{avancement-des-cinq-demandes}}

\begin{longtable}[]{@{}
  >{\raggedright\arraybackslash}p{(\columnwidth - 4\tabcolsep) * \real{0.3333}}
  >{\raggedright\arraybackslash}p{(\columnwidth - 4\tabcolsep) * \real{0.3333}}
  >{\raggedright\arraybackslash}p{(\columnwidth - 4\tabcolsep) * \real{0.3333}}@{}}
\toprule\noalign{}
\begin{minipage}[b]{\linewidth}\raggedright
Demande
\end{minipage} & \begin{minipage}[b]{\linewidth}\raggedright
Travail réalisé
\end{minipage} & \begin{minipage}[b]{\linewidth}\raggedright
Validation encore attendue
\end{minipage} \\
\midrule\noalign{}
\endhead
\bottomrule\noalign{}
\endlastfoot
Contribution scientifique & Audit technique, positionnement et cinq
questions de validation préparés & Avis de l'encadrant sur la
contribution et la nouveauté \\
Autre machine & 560 modèles recalculés, 96 contrastes reconstruits, 12
réentraînements terminés & Accepter la limite d'identité des entrées ou
demander une étude prospective plus stricte \\
Relecture et bibliographie & Relecture interne, corrections et
références actualisées & Relecture scientifique humaine indépendante \\
Revue et auteurs & Journal of Computational Science retenu comme cible
de préparation ; formulaire auteurs complet & Accord sur la revue, noms,
ordre, affiliations, contact et déclarations \\
Dossier éditorial & Manuscrit, figures, annexes, sources LaTeX,
highlights et lettre préparés & Guide complet au dépôt, intégration des
retours et accord final des auteurs \\
\end{longtable}

\hypertarget{documents-pruxeats-pour-la-relecture}{%
\section{Documents prêts pour la
relecture}\label{documents-pruxeats-pour-la-relecture}}

\begin{itemize}
\tightlist
\item
  \texttt{ARTICLE1\_MANUSCRIPT\_V11.pdf} : manuscrit anglais révisé,
  titre original conservé.
\item
  \texttt{ARTICLE1\_SUPPLEMENT\_V11.pdf} : détails de protocole et
  ensemble des 96 contrastes V10, y compris résultats non significatifs.
\item
  \texttt{SCIENTIFIC\_AUDIT.md} : corrections et points à examiner par
  l'encadrant.
\item
  \texttt{LITERATURE\_UPDATE.md} et \texttt{references.json} :
  actualisation et bibliographie modifiable.
\item
  \texttt{HIGHLIGHTS.txt} et \texttt{COVER\_LETTER\_DRAFT.md} :
  documents éditoriaux préparés.
\item
  \texttt{reproduce\_article1\_v11.py} et workflow GitHub : reproduction
  et sous-ensemble de 12 entraînements fixés avant exécution.
\end{itemize}

\hypertarget{grille-de-validation-pour-lencadrant}{%
\section{Grille de validation pour
l'encadrant}\label{grille-de-validation-pour-lencadrant}}

\begin{enumerate}
\def\labelenumi{\arabic{enumi}.}
\tightlist
\item
  Le titre et la question scientifique reflètent-ils correctement le
  projet doctoral maintenu ?
\item
  Les limites de M6 et de l'identité de correction M7 sont-elles
  correctement interprétées ?
\item
  La distinction entre résultats historiques et V09-V10 exploratoires
  est-elle suffisante ?
\item
  La contribution est-elle assez forte pour la revue retenue ? Une
  reproduction fidèle d'un concurrent ou une autre baseline de placement
  est-elle nécessaire pour les revendications maintenues ?
\item
  Les auteurs acceptent-ils le choix de la revue, leur contribution,
  leur affiliation et les déclarations ?
\end{enumerate}

Une réponse écrite à ces points permettra d'intégrer les corrections
sans transformer une relecture automatisée en validation humaine
fictive.

\hypertarget{lettre-de-soumission-pruxe9paruxe9e---uxe0-valider}{%
\section{Lettre de soumission préparée - à
valider}\label{lettre-de-soumission-pruxe9paruxe9e---uxe0-valider}}

Le brouillon ci-dessous n'est pas envoyé. Son ouverture au futur reflète
l'absence actuelle d'approbation finale. Après validation des auteurs,
l'auteur correspondant pourra adapter l'ouverture au dépôt effectif et
ajouter les déclarations confirmées.

Dear Editors of the Journal of Computational Science,

We intend to submit the research manuscript entitled ``Reference-Set
Decoupled Co-Adaptive Training for Physics-Informed Neural Networks''
for consideration as a research article, subject to final author
approval.

The study examines whether evaluating a global loss controller on a
fixed reference set improves co-adaptive PINN training. It compares this
mechanism with coupled control, proposal correction, kernel-density
volume weighting and variational residual attention on a common
computational backbone. The results do not establish universal
superiority of the fixed-reference method. Instead, the manuscript
reports configuration-dependent benefits and limitations, supported by
paired experiments, reference-solution validation and an explicit audit
of observation-loss comparability.

The experimental package includes 200 V09 ablations and 560 additional
V10 training runs. Frozen checkpoints, configurations, evaluation tables
and analysis scripts accompany the study. The manuscript distinguishes
exploratory findings, reused comparisons and computational timing
limitations. Its emphasis on controlled numerical evidence and
reproducible simulation motivates our selection of the Journal of
Computational Science.

Before submission, the corresponding author will confirm the final
author list and affiliations, originality and exclusive consideration,
funding and competing-interest declarations, and approval of the final
manuscript by all authors. The use of OpenAI's ChatGPT/Codex in code
development, analysis support and manuscript preparation is disclosed
for author verification.

Sincerely,

Youssef EL MOUTEE

Corresponding-author designation, affiliation and contact details:
pending confirmation.

\end{document}
